\section{Decomposition and exhaustion}\label{sec:decomposition-exhaustion}



Under the recognition hypothesis every well-formed decomposition record meets
the condition of Definition~\ref{def:scoped-exact-six-question}
(Theorem~\ref{thm:upper-bound-classification}), so the exact-six question
reduces to the existence of one well-formed record, which the theorem below
proves for every description. Given \(D \in \FATCD{}\), is there a
well-formed record that lists the six role channels, gives each a
status, and accounts for every remaining raw feature through an
inactive status record or a residual classification? This section shows that there is, and explains why a
channel is not the same as an active role. The lower bounds show that each
of the six channels is materially realizable in the domain; the
upper-bound classification shows that, under the recognition hypothesis, no
kind residual or record residual is licensed label (x) or (xi).

\subsection{Decomposition records and channel statuses}\label{subsec:dec-records}

\begin{definition}[Channel status]\label{def:channel-status}
For \(D\in\FATCD{}\), a decomposition record \(\Dec{D}\)
(Definition~\ref{def:decomposition-record}) and \(i \in \{1,\ldots,6\}\), the
\emph{role channel} \(\Chan{i}{D}\) is the pair of \(i\) and the status entry a
record \(\Dec{D}\) chooses for \(i\); the status value \(\Status{i}{D}\) is fixed
below from \(D\) alone, and the record only chooses the cluster, offer, seed or
declaration it names; the channel exists whether or not the status is
active. A \emph{channel status} for
\(\Chan{i}{D}\) is exactly one of the five values
\begin{itemize}[topsep=2pt,itemsep=1pt]
  \item \(\activeproj\) --- an admissible \(\pi_i(D)\) exists with \(W_i(C)\)
    and all required attachments;
  \item \(\collapsedbc\) --- no admissible \(\pi_i(D)\) exists, and an
    honestly recorded offer of a specified cluster has check result
    \(W_i(C)=\mathrm{false}\);
  \item \(\belowthr\) --- neither preceding case applies, and a covered seed for
    \(P_i\) is present;
  \item \(\outsidescope\) --- none of the preceding cases applies, and an
    honest outside-scope declaration is present;
  \item \(\absentrec\) --- none of the preceding cases applies; the
    decomposition records that absence.
\end{itemize}
In prose these are called active, collapsed, below threshold, outside scope and absent. The status name \(\collapsedbc\) refers to a failed offer; it is unrelated to
the collapse-case annotations of Definition~\ref{def:threshold-attachment} and
to the forbidden collapses of Table~\ref{tab:boundary-matrix}.
The status \(\Status{i}{D}\) depends only on \(D\); \(\Chan{i}{D}\) and
\(\Residual{D}\) are taken relative to a fixed record \(\Dec{D}\), and we write
\(\Chan{i}{\Dec{D}}\) and \(\Residual{\Dec{D}}\) when the record matters. Since
\(\Chan{i}{D}\) contains its status entry, the pair
\((\Chan{i}{D},\Status{i}{D})\) in Definition~\ref{def:decomposition-record}
repeats the status as a separate component. A \emph{covered seed} for \(P_i\) is, respectively:
\begin{enumerate}[label=(\arabic*),topsep=2pt,itemsep=1pt]
  \item a map or partial-map record whose
declared use is \emph{update}, together with a relation record of declared use quotient that it
references as its source relation, or with the relation record
serving as source relation in a cluster satisfying \(W_1\) that contains it;
  \item a predicate record
whose predicate field is declared as a macro-specification whose formula has
exactly one free variable \(r\), occurring in it and ranging over records of
the candidates' kind, and whose evaluations recorded, as in the \(W_2\) row, in a field of the realization relation or of a comparison or check record of \(D\), at declared records of
that kind, are correct (each equals the value of the formula at that record in
\(D\)) and include both truth values, together with a
realization relation (a relation record of declared use realization) whose target field is this predicate record and whose declared candidate class is a nonempty set of declared records of that kind; the seed consists of these two records, and either may be named as carrying it; the named seed fields are the formula field of the predicate record and, on
the realization-relation record, its candidate-class and target fields, its
list of satisfying members, its recorded evaluations of that formula (as datum
values or as certificates \((s(t),b)\) at declared records \(t\) of the
candidates' kind), and its certificate exhibiting a satisfying realizer or
checking that none exists; no other assertion field of either record is a named
seed field;
  \item a certificate, of one of the two types in the \(W_3\) row and naming as inputs two path or derivation records with the same declared endpoints, that these are not related by the rewrite rules recorded in \(D\), carried in a field of a content record of \(D\);
  \item a
law-bearing stage, order, transition or refinement record;
  \item a relation record whose
declared use is \emph{quotient}, or a map record whose declared use is \emph{packaging};
  \item an event, provenance or trace record with its declared event or trace
identity and reference fields, or a replay or check record with its declared
trace reference and a certificate whose formula has exactly one of the two forms prescribed for the \(W_6\) row after the witness table, with that trace as \(\tau\) (annotations are not
seeds).
\end{enumerate}
Unlike the no-realizer branch of the \(W_2\) row, clause (2) does not require
the candidate class to be the full declared domain of a packaging or update
map. Every closed check-language formula without \(W_j\) atoms admits the
seed-(2) presentation of Remark~\ref{rem:recognition-open}, provided the completed wrapped description lies in \FATCD{} and contains two distinct
declared records of one kind (paragraph \emph{Wrapping a closed formula as a \Ptwo{} seed}). Such a seed
supplies recognized \Ptwo{} content but does not by itself certify \(W_2(C)\) or
an active \Ptwo{} projection. In ``covered seed'', ``covered''
names the seed clauses (1)--(6); it is not a reference to the covered scope.
\par The named seed fields are: in (1), the graph of the update map, any
source-relation reference it carries, and the class data of the source
relation; in (3), the
fields of the comparison certificate; in (4), the law certificate and the
endpoint or index fields it checks; in (5), the class data of the relation or
the graph and quotient-object reference of the map; and in (2) and (6), the
fields named there.
\par A record \emph{carries} a seed when it is the map record in
(1), the predicate or realization-relation record in (2), the record holding
the certificate in (3), the law-bearing record in (4), the relation or map
record in (5), or the event, provenance, trace, replay or check record in (6);
a named seed field lying on another record is an input to that record's
seed. A comparison or check record whose recorded evaluations of a clause-(2) formula are used to meet that clause does not carry that seed. This does not exclude recognition through other fields or role data. Its classification must be checked against all the criteria of Definition~\ref{def:residual-scheme}; if none of (i)--(ix) holds, it is licensed (x). By the witness table, every cluster satisfying \(W_i\) contains a
covered seed for \(P_i\) (for \(i=2\) the row, like seed clause (2), requires the recorded evaluations
to be correct): the update map with its source relation data, the
macro-specification with its realization relation, the comparison certificate
of the two routes, the law-bearing record, the quotient relation or packaging
map, and the event or trace record, respectively. A \emph{failed offer} is an
honestly recorded attempt to use a specified cluster \(C\) as a \(P_i\)
witness for which the recorded check gives \(W_i(C)=\mathrm{false}\). An \emph{outside-scope declaration} for \(P_i\) is a claim annotation naming
\(P_i\) and a record \(r\), with no witness-cluster field (so it is not an offer
and its audit records no check of \(W_i\)), whose referenced record \(r\) lies outside the covered scope (the
criterion after Definition~\ref{def:residual-scheme}) and whose kind occurs in
the \(W_i\) row of the witness table; its honest audit lists the
outside-scope field of \(r\) in its non-certification field. The status of \(\Chan{i}{D}\) is the
first that applies in the order: \(\activeproj\) if some cluster supports
an admissible derived projection \(\pi_i(D)\) in the sense of
Definition~\ref{def:role-projection}, including \(W_i(C)\) and every
applicable attachment of Definition~\ref{def:threshold-attachment}; \(\collapsedbc\) if there is a
failed offer; \(\belowthr\) if a covered seed is present;
\(\outsidescope\) if there is an outside-scope declaration; and
\(\absentrec\) otherwise. An inactive status is recorded in \(\Dec{D}\); it
is not a role witness. We write \(\Status{i}{D}\) for the channel status of
\(\Chan{i}{D}\).
\end{definition}

Statuses are assigned only by the ordered rule above. An admissible
projection takes priority over every failed offer or incomplete seed. In its
absence, a failed offer gives \(\collapsedbc\); otherwise a covered seed
lacking a complete admissible projection gives \(\belowthr\). Only then are an
outside-scope declaration and the absent fallback considered.
Proposition~\ref{prop:boundary-distinctions} uses the same two labels for
local attempt verdicts, which need not equal channel statuses. In the running
example, the \Pfour{} channel is \(\belowthr\): its transition record is
present, but no threshold annotation targets a cluster containing it.

\begin{definition}[Decomposition/exhaustion record]\label{def:decomposition-record}
A \emph{decomposition/exhaustion record} for \(D \in \FATCD{}\) is a tuple
\[
\Dec{D} = \bigl(\Raw{D},\; (\Chan{i}{D}, \Status{i}{D})_{i=1}^{6},\;
\Residual{D},\; \mathrm{Classify},\; \mathcal{L},\; \mathcal{V},\;
\Audit{D},\; \mathcal{N}\bigr)
\]
where
\begin{itemize}[topsep=2pt,itemsep=1pt]
  \item the six role channels together with their statuses form the
    six-channel status sequence, each \(\Status{i}{D}\) being the value that the
    ordered rule of Definition~\ref{def:channel-status} assigns. If \(\Status{i}{D} = \activeproj\), the channel
    carries an active projection \(\pi_i(D)\) with its full attachment
    data. If \(\Status{i}{D}\) is inactive, the channel carries the
    corresponding below-threshold, collapsed-boundary,
    absent-with-record, or outside-scope status record;
  \item \(\Residual{D} \subseteq \Raw{D}\) is the residual feature set,
    defined as follows. Let \(A(\Dec{D})\) be the union of the content
    clusters of the active projections and the attachment annotations (the
    threshold, witness-schema, level and collapse-case annotations attached to
    the cluster, the claim annotation naming \(P_i\) and the cluster, the audit
    annotations referring to that claim, the nonclaim annotations targeting it, the translation and lift or
    forgetting annotations that Definition~\ref{def:role-projection} requires for
    a level crossing of the projection, and,
    for \(P_6\), the event-or-trace claim and honest audit used to establish
    the \(W_6\) audit-entry conjunct)
    that the channel entry assigns to its projection (an active entry assigns at least one annotation of each kind that Definition~\ref{def:role-projection} requires for admissibility and may leave further annotations targeting the cluster unassigned). Let \(S(\Dec{D})\) be the union of the
    clusters and offers explicitly named by the inactive status entries: the
    records of the failed offered cluster with the offer's claim annotation and the audits referring to it, for \(\collapsedbc\); all records of the covered seed it names (for clauses (1) and (2), both records the clause names), for
    \(\belowthr\); and the outside-scope declaration with the audits referring to it, for
    \(\outsidescope\); \(\absentrec\) accounts for no feature, and a status
    entry accounts only for records it names. Then
    \(\Residual{D}=\Raw{D}\setminus\bigl(A(\Dec{D})\cup S(\Dec{D})\bigr)\),
    relative to the chosen record, and \(\mathrm{Classify}\) assigns to each
    \(r \in \Residual{D}\), including any bookkeeping annotation not assigned
    to a channel, a label in the scheme of
    Definition~\ref{def:residual-scheme}; already covered bookkeeping may
    receive (ix);
  \item \(\mathcal{L}\) collects the loss and lift records for any level
    crossings involved in projections or classifications;
  \item \(\mathcal{V}\) collects instrument-visibility and
    empirical-bridge annotations where these are asserted;
  \item \(\Audit{D}\) is the honest bookkeeping subrecord;
  \item \(\mathcal{N}\) is the explicit nonclaim record.
\end{itemize}
A decomposition/exhaustion record is \emph{well-formed} when each
\(\Status{i}{D}\) is drawn from the five values of
Definition~\ref{def:channel-status}, each active
\(\pi_i(D)\) is admissible in the sense of
Definitions~\ref{def:role-projection} and~\ref{def:threshold-attachment},
each inactive status record names a failed offer, covered seed or
outside-scope declaration of \(D\) for that channel, as
Definition~\ref{def:channel-status} requires (none for \(\absentrec\)), and every raw feature \(r\) lies in \(A(\Dec{D})\), or in \(S(\Dec{D})\), or in
\(\Residual{D}\) and receives a label, and, for each
\(r\in\Residual{D}\), \(\mathrm{Classify}(r)\) is a label whose criterion in the
table after Definition~\ref{def:residual-scheme} holds for \(r\) relative to
\(\Dec{D}\). The records in
\(\mathcal{L}\) and \(\mathcal{V}\) are auxiliary annotations on this
coverage, not additional role channels. The status entries and the
collections \(\mathcal{L}\), \(\mathcal{V}\) and \(\mathcal{N}\) are data
of \(\Dec{D}\). Each cluster, seed, offer or declaration named by a status is
selected from, or supported by, the records of \(\Raw{D}\) and \(\Audit{D}\)
required above; the selection and the status entries need not themselves be
raw records. Forming \(\Dec{D}\) does not change \(D\).
\end{definition}

By the ordered rule of Definition~\ref{def:channel-status}, the status
\(\Status{i}{D}\) depends only on \(D\), not on the decomposition record; the
cluster, offer, seed or declaration that a record names, and hence
\(A(\Dec{D})\), \(S(\Dec{D})\) and \(\Residual{D}\), may depend on the record.
Because \(\Residual{D}\) leaves out every feature already accounted for by
the active projection named for its channel or by an inactive status record, labels~(i)
and~(ii) of Definition~\ref{def:residual-scheme} do not occur for a residual
feature of \(\Dec{D}\): both require the record to lie in a cluster named for
an active projection. This is why the results below list labels (iii)--(ix).

\paragraph{Two residual sets.} A residual feature in the sense of
Definition~\ref{def:residual-scheme} is also called a \emph{kind residual}, in
statements and elsewhere, and an element of \(\Residual{D}\) a \emph{record
residual}.
Membership in the first depends on the raw-grammar kind of a feature, and
membership in the second on whether the record \(\Dec{D}\) accounts for it
through a channel. In particular, a bookkeeping annotation is not a kind
residual but may be a record residual if the chosen channel assignments do not
account for it. A kind residual absorbed into an active projection is not a
record residual, which is why label~(i) does not occur among record residuals.
The recognition hypothesis and
Theorem~\ref{thm:upper-bound-classification} cover both sets.


\subsection{Existence, six-channel exactness, and residual exhaustion}\label{subsec:dec-theorems}

\begin{theorem}[Decomposition existence]\label{thm:decomposition-existence}
Every description \(D \in \FATCD{}\) admits a well-formed decomposition/exhaustion
record \(\Dec{D}\) in the sense of
Definition~\ref{def:decomposition-record}.
\end{theorem}

\begin{proof}
For each \(i \in \{1,\ldots,6\}\), range over the finitely many clusters
\(C\subseteq\Raw{D}\). If some \(C\) satisfies \(W_i(C)\) together with every
item of Definitions~\ref{def:threshold-attachment}
and~\ref{def:role-projection}, name one such \(C\), record
\(\Status{i}{D} = \activeproj\) and the active projection
\(\pi_i(D)\) it determines; otherwise record the inactive status that the
ordered rule of Definition~\ref{def:channel-status} assigns, and name one
failed offer, covered seed or outside-scope declaration for the first case of
the rule that applies. Six status entries
result, one per role. The raw features absorbed by these channel
records are removed from the residual stage: the statuses, their
supporting clusters and their assigned annotations are selected first, and
\(\Residual{D}\) is then formed by the set difference of
Definition~\ref{def:decomposition-record}. For each \(r\) in it, choose a
licensed label among (i)--(ix) if one applies. If none applies, then \(r\) is
not outside the covered scope (otherwise (viii) would apply), \(D\) is finite,
and every claim concerning \(r\) is referred to by an honest audit annotation; so the criterion of
(x) holds. Thus \(\mathrm{Classify}\) is total, and the record is well-formed
irrespective of which labels occur. An unassigned bookkeeping annotation is
assessed by the same criteria, ordinarily (ix). Relative to the chosen
record, (i) cannot apply to a record residual, since absorbed content was
placed in \(A(\Dec{D})\); under the recognition hypothesis
(Definition~\ref{def:recognition-hypothesis}),
Theorem~\ref{thm:upper-bound-classification} further places every label in
(iii)--(ix), with (x) and (xi) empty. The records \(\mathcal{L}\) and \(\mathcal{V}\) are
read from \(\Raw{D}\) and \(\Audit{D}\); the bookkeeping subrecord is
\(\Audit{D}\); and \(\mathcal{N}\) records the nonclaims carried
forward from the witness discipline. The resulting tuple is a
decomposition/exhaustion record, well-formed by construction.
Non-vacuity of the six channels follows from
Theorem~\ref{thm:six-role-lower-bounds}, which supplies, for each
channel, a description in which that channel is active.
\end{proof}

\begin{theorem}[Six-channel exactness and residual exhaustion]\label{thm:six-channel-exactness}
For every \(D \in \FATCD{}\) and every well-formed
\(\Dec{D}\) in the sense of
Definition~\ref{def:decomposition-record}:
\begin{enumerate}[label=(\roman*),topsep=2pt,itemsep=1pt]
  \item the channel sequence \((\Chan{i}{D})_{i=1}^{6}\) has length
    exactly six, one per primitive role \Pone{}, \ldots, \Psix{};
  \item every raw feature of \(D\) is accounted for at the
    channel/exhaustion level by an active projection \(\pi_i(D)\), a
    channel-status record for an inactive channel, or a residual
    classification under the scheme of
    Definition~\ref{def:residual-scheme}, possibly with label (x) or (xi);
    since label (x) applies to in-scope features for which no criterion
    (i)--(ix) holds, this clause alone does not exclude an unresolved
    candidate (see (iii)); the records in \(\mathcal{L}\)
    and \(\mathcal{V}\) are auxiliary annotations on this coverage;
  \item under the recognition hypothesis
    (Definition~\ref{def:recognition-hypothesis}) and within the covered scope,
    no raw feature is classified as an unresolved seventh-role threat or as
    passing the recorded seventh-role screen
    (Theorem~\ref{thm:upper-bound-classification}); the three principal threat
    classes are closed without the recognition hypothesis, in the
    recognized forms of
    Propositions~\ref{prop:probability-closure}--\ref{prop:agency-closure};
  \item the status sequence \((\Status{i}{D})_{i=1}^{6}\) is the same for every
    well-formed \(\Dec{D}\).
\end{enumerate}
\end{theorem}

\begin{proof}
Exactness of the channel sequence is built into the definition of
\(\Dec{D}\); there are six role channels by
Definition~\ref{def:decomposition-record}. Coverage of raw features is
the well-formedness condition of
Definition~\ref{def:decomposition-record} together with totality of
\(\mathrm{Classify}\) on \(\Residual{D}\), which holds by the case argument in
the proof of Theorem~\ref{thm:decomposition-existence}. Emptiness of categories (x) and (xi),
under the recognition hypothesis (Definition~\ref{def:recognition-hypothesis})
and within the covered scope, is
Theorem~\ref{thm:upper-bound-classification} directly. The last assertion of
(iii) is Propositions~\ref{prop:probability-closure},
\ref{prop:causality-closure} and~\ref{prop:agency-closure}, which hold for
every \(D\in\FATCD{}\) and every well-formed \(\Dec{D}\) without the
recognition hypothesis. For (iv), Definition~\ref{def:decomposition-record}
records each status by the ordered rule of Definition~\ref{def:channel-status},
whose clauses refer only to \(D\).
\end{proof}

\subsection{Channel-versus-activation discipline}\label{subsec:channel-vs-activation}

Theorem~\ref{thm:six-channel-exactness} asserts six channels in every
decomposition/exhaustion record. It does not assert six active
projections in every record, nor does it assert that a well-formed
record is unique. The distinction between channel presence and
channel activation prevents six-channel exactness from being read as
an all-six-active claim.

Membership of \(D_1,\dots,D_6\) in \(\FATCD_{\mathrm{RH}}\) is verified record by
record in the proof of Theorem~\ref{thm:scoped-exact-six}; that verification
uses only Theorem~\ref{thm:six-role-lower-bounds} and the definitions. In brief:
every content record of \(D_i\) contributes a required field to \(C_i\) or
carries a covered seed, each of its assertion fields being required or a named
seed field (R1), or it has no assertion field (R1 or R3); every structural
annotation of \(D_i\) carries only its typing, target and label or level field,
and every bookkeeping annotation only datum and reference fields besides its
exempt fields (R2).

\begin{proposition}[Channel-versus-activation discipline]\label{prop:channel-vs-activation}
Let \(D \in \FATCD{}\) and let \(\Dec{D}\) be a well-formed decomposition/exhaustion
record. The following are distinct statements:
\begin{enumerate}[label=(\roman*),topsep=2pt,itemsep=1pt]
  \item \emph{Six channels}: the channel sequence
    \((\Chan{i}{D})_{i=1}^{6}\) has length six with the fixed role names
    \Pone{}, \ldots, \Psix{}.
  \item \emph{Six active projections}: for each \(i\),
    \(\Status{i}{D} = \activeproj\) and the channel carries an active
    \(\pi_i(D)\).
\end{enumerate}
Theorem~\ref{thm:six-channel-exactness} establishes (i) for every
\(D \in \FATCD{}\). It does not establish (ii). In particular, an
admissible \(D \in \FATCD{}\) may have any number of active channels
between zero and six; inactive channels are recorded with status
\(\belowthr\), \(\collapsedbc\), \(\absentrec\), or \(\outsidescope\)
and are not active witnesses for the corresponding roles. For each
\(k\in\{0,\dots,6\}\) some \(D\in\FATCD_{\mathrm{RH}}\) has exactly \(k\) active
channels, and some single description in \(\FATCD_{\mathrm{RH}}\) has channels of all
five statuses.
\end{proposition}

\begin{proof}
Statements (i) and (ii) differ whenever some channel is inactive, which the
last clause provides for every \(k<6\). \par\emph{Exactly \(k\) active channels.} For \(k=0\), let \(D_0\) consist of a
single object record: it is finite and closed, it has no claims, so it is
audited, and it has no cluster satisfying any \(W_i\), no failed offer, no
covered seed and no outside-scope declaration, so the status rule of
Definition~\ref{def:channel-status} gives every channel \(\absentrec\). Its
object record has no witness field and no predicate, comparison, invariant,
law or claim field, so it satisfies (R3) of
Definition~\ref{def:recognition-hypothesis} and
\(D_0\in\FATCD_{\mathrm{RH}}\). For
\(1\le k\le 6\), choose \(k\) distinct indices \(S\) and let \(D_S\) be the
disjoint union of the descriptions \(D_i\), \(i\in S\), of
Theorem~\ref{thm:six-role-lower-bounds}, with record identifiers made
disjoint by an injective renaming applied recursively to every
record-identifier occurrence in reference, datum, term-valued and
formula-valued fields, including identifiers inside finite lists,
tuples, relation classes and map tables, and the argument lists of
\(W_j\) atoms. Rational constants and other non-identifier values
are unchanged. A numeral in record-sorted position (paragraph \emph{Sorts} of Section~\ref{sec:primitive-role-architecture}), such as \(0\) and \(2\) in the formula of \(s\) in \(D_2\), is a record-identifier occurrence and is renamed; a numeral in rational-sorted position, such as an element of the monoid list of \(D_3\), is not. Each condition of Definition~\ref{def:fatcd} concerns a record, an
annotation or a claim together with the records it refers to, which lie in one component, except that certificate correctness and the evaluation of \(W_j\) atoms may read description-wide data (all declared records or those of one kind, the rule list \(R\), \(\mathrm{Audit}_0(D)\)); the next sentences check that this data leaves every recomputed value unchanged. The certificates' quantifiers are bounded by listed object records or finite field lists, except \(\forall u\in|E|\), which is guarded by \(E(u,u)\); this is false at every record of another component, since renaming keeps the classes of \(E\) inside its component. So the content checks retain their values in \(D_S\). The \(W_1,W_2,W_4,W_5\) checks remain within their components. The
\(W_3\) row also requires the rules of its cluster to be all rewrite rules
recorded in the description; since the indices in \(S\) are distinct and \(D_3\)
is the only component recording rewrite rules, this clause keeps its value in
\(D_S\). For
\(W_6(C_6)\), the honest event-or-trace audit entry of \(D_6\) remains in
\(\mathrm{Audit}_0(D_S)\), so its separate audit conjunct also remains true.
Thus the audits remain honest and \(D_S\in\FATCD{}\). For \(i\in S\), the cluster \(C_i\)
with its attachments still satisfies \(W_i\), so channel \(P_i\) is
\activeproj{}. For \(j\notin S\), no component has a witness-schema
annotation naming \(P_j\), together with a threshold and an honest audit of
\(W_j(C)\) for a cluster \(C\), and making record identifiers disjoint creates
no such annotation or cross-component reference. So no admissible
\(\pi_j(D_S,C)\) in the sense of Definition~\ref{def:role-projection} exists,
and channel \(P_j\) is not active. Hence \(D_S\) has exactly \(k\) active
channels.
Each of these descriptions lies in \(\FATCD_{\mathrm{RH}}\).
The inventory in the proof of the last clause of
Theorem~\ref{thm:scoped-exact-six} (Section~\ref{sec:scoped-exact-six}, p.~\pageref{thm:scoped-exact-six}), which uses
Theorem~\ref{thm:six-role-lower-bounds} and the definitions but not this
proposition, establishes (R1) or (R3) for every content record of each \(D_i\)
using records and checks within that component, and (R2) for its
structural and bookkeeping annotations. Under this renaming each
certificate, claim statement and audit conclusion of \(D_i\) becomes its renamed
image, which recomputes to the same value in \(D_S\). They apply to every
kind residual and every possible record residual of \(D_S\),
independently of the chosen decomposition record. The case \(D_0\)
was checked above.
\par\emph{All five statuses.} For the five statuses, \(D_S\) supplies \(\activeproj\) and \(D_0\) supplies
\(\absentrec\). Let \(o\) be the object record of \(D_0\). Let \(D_{\mathrm c}\) consist of state
records \(x\ne y\), an object record \(X_{\mathrm c}\) listing \(x\) and \(y\), a
relation record \(E\) on \(X_{\mathrm c}\) of declared use quotient with the
single class \(\{x,y\}\), a claim naming \Pfive{} and \(C=\{x,y,E\}\), and an
honest audit listing \((W_5(C),\mathrm{false})\); no packaging map is present, so
this failed offer lies on a cluster carrying a covered \Pfive{} seed, and channel
\Pfive{} is \(\collapsedbc\). Here \(x\), \(y\) and \(X_{\mathrm c}\) satisfy (R3), \(E\) carries a
covered seed whose class data are named seed fields and satisfies (R1), and the
annotations satisfy (R2). Let \(D_{\mathrm{o}}\) be \(D_0\) together with a predicate record \(Q\) (here \(Q\) is a predicate record, unrelated to the quotient objects \(Q\) of the \(W_5\) row and Remark~\ref{rem:recognition-open}, and \(o\) is the object record of \(D_0\), not the order record of \(D_{\mathrm{scr}}\)) of arity
one, its argument place declared to range over object records, with no formula field and no relation record presenting it, a predicate
record \(s\) whose only assertion field is \(Q(o)\), a claim naming \Ptwo{} and
\(s\) with statement \(Q(o)\), and an honest audit that lists no certificate,
records no check of \(W_2\) (so the claim is not an offer), concludes \(o=o\)
and lists the field of \(s\) in its non-certification field. This claim is an
outside-scope declaration, so channel \Ptwo{} of \(D_{\mathrm{o}}\) is
\(\outsidescope\). The declaration is not an admissible claim in the sense of
Definition~\ref{prop:honest-bookkeeping}, since no audit certifies \(Q(o)\);
Definition~\ref{def:channel-status} requires only that its audit be honest, and
the conclusion \(o=o\) makes the audit honest while certifying nothing about
\(Q\). In \(D_{\mathrm{o}}\) \(o\) and \(Q\)
satisfy (R3), \(s\) satisfies (R5), and the annotations satisfy (R2). The
description \(D_{\mathrm{ex}}\) of Section~\ref{sec:fatcd-domain} supplies
\(\belowthr\) for its \Pfour{} channel. It lies in \(\FATCD_{\mathrm{RH}}\). The records \(F\), \(E\), the
split-pair comparison, \(q\), \(\{A,B\}\), the quotient-validity check, the
event, the trace and the check record each contribute a required field to
\(C_1\), \(C_5\) or \(C_6\), and their assertion fields (the split-pair,
validity and replay certificates) are required for \(W_1\), \(W_5\) and \(W_6\)
respectively, so these records satisfy (R1). The provenance record carries a
covered \Psix{} seed, and the transition record carries a covered \Pfour{} seed
whose law field is a named seed field; both satisfy (R1). The state records,
\(S_0\), the object records \(A\), \(B\) and the replay record have no assertion field and satisfy (R1) or
(R3); the index record satisfies (R4) by the one-step chain to \(E\), which satisfies
(R1); and its annotations satisfy (R2).
\par\emph{All five in one description.} Let \(D_{\mathrm{all}}\) be the renamed disjoint union of \(D_6\),
\(D_{\mathrm c}\), \(D_{\mathrm o}\) and \(D_T\), where \(D_T\) has a state record
\(z\), an object record \(Z\) listing it, and a transition record \(T:Z\to Z\),
\(T(z)=z\), whose law field carries the certificate, recomputed true, that \(T\)
is defined on \(Z\) with values in \(Z\). The argument for \(D_S\) gives
\(D_{\mathrm{all}}\in\FATCD{}\); \(T\) carries a covered \Pfour{} seed with named
law field and satisfies (R1), \(z\) and \(Z\) satisfy (R1) or (R3), and the other
components are inventoried above, so \(D_{\mathrm{all}}\in\FATCD_{\mathrm{RH}}\).
No component supplies a \Pone{}, \Ptwo{} or \Pthree{} seed, an active projection
other than that of \(D_6\), or a failed offer other than that of
\(D_{\mathrm c}\); so by the ordered rule \Psix{} is \(\activeproj\), \Pfive{}
\(\collapsedbc\), \Pfour{} \(\belowthr\), \Ptwo{} \(\outsidescope\), and \Pone{}
and \Pthree{} \(\absentrec\).
\end{proof}

\begin{proposition}[Decomposition records are not unique]\label{prop:non-uniqueness}
A description \(D \in \FATCD{}\) may admit more than one well-formed
decomposition/exhaustion record in the sense of
Definition~\ref{def:decomposition-record}. In particular, distinct
well-formed records may name different witness clusters for one active
channel, and then have different sets \(A(\Dec{D})\), \(\Residual{D}\) and
residual labels; Definition~\ref{def:decomposition-record} also leaves the choice among
level annotations attached to the named cluster and the placement of
loss/lift and visibility records to the record. Such a description can be chosen in \(\FATCD_{\mathrm{RH}}\).
Theorems~\ref{thm:decomposition-existence}
and~\ref{thm:six-channel-exactness} assert existence, not uniqueness.
\end{proposition}

\begin{proof}
Let \(D\) be the disjoint union of two copies of the description \(D_5\) of
Theorem~\ref{thm:six-role-lower-bounds} (with the same renaming), as in the proof of
Proposition~\ref{prop:channel-vs-activation}; then \(D\in\FATCD{}\), and each
copy carries a cluster, \(C_5\) and \(C_5'\), that satisfies \(W_5\) with its
attachments. The construction of Theorem~\ref{thm:decomposition-existence}
may name either cluster for the \Pfive{} channel. In the record naming
\(C_5\), the content records of \(C_5'\) that contribute a required field to
\(C_5'\), a six-role-aligned cluster that the chosen projection does not
absorb, receive label (v), and any others, having no predicate, comparison,
invariant, law or claim field and supplying no witness field, receive (ix-c); in the record naming \(C_5'\), the same holds with the
copies exchanged. The \(W_5\) assertion of the unchosen copy is represented by
its check record and other required-field records, labelled (v); its
bookkeeping annotations receive (ix-a) and its threshold and level annotations
(ix-b).
The two records name different clusters, so they differ, and both are
well-formed. The description \(D\) satisfies the recognition hypothesis: each
content record of either copy supplies a witness field to that copy's cluster
satisfying \(W_5\), or is an input without assertion fields, and so satisfies
(R1) or (R3) as in the inventory in the proof of
Theorem~\ref{thm:scoped-exact-six} (p.~\pageref{thm:scoped-exact-six}); its annotations satisfy (R2). Hence
\(D\in\FATCD_{\mathrm{RH}}\).
\end{proof}

\paragraph{The running example, decomposed.} For the description
\(D_{\mathrm{ex}}\) of Section~\ref{sec:fatcd-domain}, a well-formed
decomposition record assigns the six channels the following statuses; the
records named are those of the table in Section~\ref{sec:fatcd-domain}. The
channels \Pone{}, \Pfive{} and \Psix{} have status \activeproj{}: the split pair
\((0,1)\) witnesses the descent obstruction, the quotient object \(\{A,B\}\)
with its packaging map witnesses the packaging, and the replayed and checked
trace witnesses the audit. The channel \Pfour{} has status \belowthr{}: the
transition record of \(F\) is present, but no threshold annotation targets a
cluster containing it. The channels \Ptwo{} and \Pthree{} have status \absentrec{}: the
description records no macro-specification and no path or derivation record. The index
record copying the classes of \(E\) is a residual feature, classified as a
presentation variant of the relation record. So \(D_{\mathrm{ex}}\) carries six
channels and three active projections, as Proposition~\ref{prop:channel-vs-activation}
allows.

\begin{remark}[What the decomposition theorems do and do not prove]\label{rem:dec-scope}
Theorems~\ref{thm:decomposition-existence}
and~\ref{thm:six-channel-exactness} establish the structural side of
the scoped exact-six claim: every \(D \in \FATCD{}\) admits a
six-channel decomposition/exhaustion record, with residual exhaustion in the
covered scope under the recognition hypothesis
(Definition~\ref{def:recognition-hypothesis}). Together with the lower-bound theorem
(Section~\ref{sec:lower-bounds}) and the upper-bound classification
(Section~\ref{sec:upper-bound-classification}), they prepare the
scoped exact-six theorem but do not replace the integrated-stress
check that follows. They do not assert that every description
activates all six roles, that the decomposition is canonical or
unique, or that the scoped exact-six program extends beyond \FATCD{};
those remain nonclaims.
\end{remark}

